\documentclass[10pt,a4paper]{amsart}
\usepackage{amsmath,amssymb,amsthm,graphicx}
% Retained preamble for the Stacks Project chapter morphisms.tex (Morphisms of Schemes).
% The chapter file pinned-original-source.tex carries only \input{preamble} on its first line, so no
% preamble text lives inside the chapter itself.  The Stacks Project uses one single shared preamble for
% every chapter; the snapshot stages it verbatim in sources/stacks-chow/preamble.tex.  The block below is
% copied mechanically, line for line, from that file's lines 14-18 (the Xy-pic setup) and lines 154-217
% (the native theorem declarations together with the mathematical macros).  Nothing was added, removed or
% redefined.  The wrapper class and the amsmath/amssymb/amsthm/graphicx base packages are supplied by the
% assembler, so the class/\externaldocument part of the shared preamble is not repeated here.
% For commutative diagrams we use Xy-pic
\usepackage[all]{xy}

% We use 2cell for 2-commutative diagrams.
\xyoption{2cell}
% Theorem environments.
%
\theoremstyle{plain}
\newtheorem{theorem}[subsection]{Theorem}
\newtheorem{proposition}[subsection]{Proposition}
\newtheorem{lemma}[subsection]{Lemma}

\theoremstyle{definition}
\newtheorem{definition}[subsection]{Definition}
\newtheorem{example}[subsection]{Example}
\newtheorem{exercise}[subsection]{Exercise}
\newtheorem{situation}[subsection]{Situation}

\theoremstyle{remark}
\newtheorem{remark}[subsection]{Remark}
\newtheorem{remarks}[subsection]{Remarks}

\numberwithin{equation}{subsection}

% Macros
%
\def\lim{\mathop{\mathrm{lim}}\nolimits}
\def\colim{\mathop{\mathrm{colim}}\nolimits}
\def\Spec{\mathop{\mathrm{Spec}}}
\def\Hom{\mathop{\mathrm{Hom}}\nolimits}
\def\Ext{\mathop{\mathrm{Ext}}\nolimits}
\def\SheafHom{\mathop{\mathcal{H}\!\mathit{om}}\nolimits}
\def\SheafExt{\mathop{\mathcal{E}\!\mathit{xt}}\nolimits}
\def\Sch{\mathit{Sch}}
\def\Mor{\mathop{\mathrm{Mor}}\nolimits}
\def\Ob{\mathop{\mathrm{Ob}}\nolimits}
\def\Sh{\mathop{\mathit{Sh}}\nolimits}
\def\NL{\mathop{N\!L}\nolimits}
\def\CH{\mathop{\mathrm{CH}}\nolimits}
\def\proetale{{pro\text{-}\acute{e}tale}}
\def\etale{{\acute{e}tale}}
\def\QCoh{\mathit{QCoh}}
\def\Ker{\mathop{\mathrm{Ker}}}
\def\Im{\mathop{\mathrm{Im}}}
\def\Coker{\mathop{\mathrm{Coker}}}
\def\Coim{\mathop{\mathrm{Coim}}}

% Boxtimes
%
\DeclareMathSymbol{\boxtimes}{\mathbin}{AMSa}{"02}

%
% Macros for moduli stacks/spaces
%
\def\QCohstack{\mathcal{QC}\!\mathit{oh}}
\def\Cohstack{\mathcal{C}\!\mathit{oh}}
\def\Spacesstack{\mathcal{S}\!\mathit{paces}}
\def\Quotfunctor{\mathrm{Quot}}
\def\Hilbfunctor{\mathrm{Hilb}}
\def\Curvesstack{\mathcal{C}\!\mathit{urves}}
\def\Polarizedstack{\mathcal{P}\!\mathit{olarized}}
\def\Complexesstack{\mathcal{C}\!\mathit{omplexes}}
% \Pic is the operator that assigns to X its picard group, usage \Pic(X)
% \Picardstack_{X/B} denotes the Picard stack of X over B
% \Picardfunctor_{X/B} denotes the Picard functor of X over B
\def\Pic{\mathop{\mathrm{Pic}}\nolimits}
\def\Picardstack{\mathcal{P}\!\mathit{ic}}
\def\Picardfunctor{\mathrm{Pic}}
\def\Deformationcategory{\mathcal{D}\!\mathit{ef}}

% --- appended 2026-10-12 (Chief): the project own definitions of the reference/slogan/history
% display environments, copied mechanically from the shared Stacks preamble lines 8-12
% (\usepackage{verbatim} plus the three \newenvironment lines, which the chapter preamble keeps
% verbatim from the project preamble). Nothing added, removed or redefined: these environments
% carry no mathematical content and the project itself typesets them as comments. Without these
% three lines a counted statement whose own text uses them fails with Environment slogan undefined.
\usepackage{verbatim}
\newenvironment{reference}{\comment}{\endcomment}
\newenvironment{slogan}{\comment}{\endcomment}
\newenvironment{history}{\comment}{\endcomment}
\usepackage{cleveref}
\crefname{originallemma}{original lemma}{original lemmas}
\Crefname{originallemma}{Original lemma}{Original lemmas}
\makeatletter
\newlabel{lemma-affine-separated}{{01S7}{source}{}{sourcepointer-1}{}}
\newlabel{lemma-affine-separated@cref}{{[originallemma][0][]01S7}{[1][0][]source}{}{}{}}
\makeatother
\crefname{originallemma}{original lemma}{original lemmas}
\Crefname{originallemma}{Original lemma}{Original lemmas}
\makeatletter
\newlabel{lemma-composition-universal-homeomorphism}{{0CEV}{source}{}{sourcepointer-2}{}}
\newlabel{lemma-composition-universal-homeomorphism@cref}{{[originallemma][0][]0CEV}{[1][0][]source}{}{}{}}
\makeatother
\crefname{originallemma}{original lemma}{original lemmas}
\Crefname{originallemma}{Original lemma}{Original lemmas}
\makeatletter
\newlabel{lemma-homeomorphism-affine}{{04DE}{source}{}{sourcepointer-3}{}}
\newlabel{lemma-homeomorphism-affine@cref}{{[originallemma][0][]04DE}{[1][0][]source}{}{}{}}
\makeatother
\crefname{originallemma}{original lemma}{original lemmas}
\Crefname{originallemma}{Original lemma}{Original lemmas}
\makeatletter
\newlabel{lemma-universal-homeomorphism}{{04DF}{source}{}{sourcepointer-4}{}}
\newlabel{lemma-universal-homeomorphism@cref}{{[originallemma][0][]04DF}{[1][0][]source}{}{}{}}
\makeatother
\crefname{originallemma}{original lemma}{original lemmas}
\Crefname{originallemma}{Original lemma}{Original lemmas}
\makeatletter
\newlabel{schemes-lemma-section-immersion}{{01KT}{source}{}{sourcepointer-5}{}}
\newlabel{schemes-lemma-section-immersion@cref}{{[originallemma][0][]01KT}{[1][0][]source}{}{}{}}
\makeatother
\title{In a composite of two morphisms, if two of the three are universal homeomorphisms then so is the third}
\author{The Stacks Project Authors}
\date{stacks-project:morphisms.tex:local-snapshot-20261012; GNU FDL 1.2 or later}
\begin{document}
\maketitle
\noindent\textit{Editorial note.} All delivered mathematical text is the author's own, copied from the pinned Stacks Project chapter file pinned-original-source.tex; no statement or proof step was rewritten, reordered or removed.\par\medskip
\noindent\textit{Editorial note.} Following the multi-claim convention for this source, source\_url is the source-level URL https://stacks.math.columbia.edu/ and the source label lemma-permanence-universal-homeomorphism is carried in claim\_key and here rather than in a per-claim tag URL.\par\medskip
\noindent\textit{Editorial note.} The preamble is the chapter's shared Stacks Project preamble, retained as sources/stacks-morphisms/retained-preamble.tex: a line-for-line excerpt of the shared Stacks preamble staged with this snapshot. No macro was invented or redefined, and this card adds no preamble text of its own.\par\medskip
\noindent\textit{Editorial note.} Statement span 11433-11439 and proof span 11441-11488 come from the halfA pool card operations/pool-cards/stacks-morphisms-halfA.json and were re-read in the pinned chapter file before this card was written.\par\medskip
\noindent\textit{Editorial note.} The proof contains two large xymatrix diagrams of base changes; both are copied verbatim and the retained chapter preamble loads Xy-pic.\par\medskip
\noindent\textit{Editorial note.} This card was previously skipped because the shared gate substring-matched my chapter name against the published items of more-morphisms.tex; the gate was fixed to exact file-name equality by the Lead and the candidate was re-gated (RESULT=OK) before being opened.\par\medskip
\noindent\textit{Editorial note.} 5 cross-references occur inside the counted spans; they are listed in external\_references as navigation pointers only.\par\medskip
\noindent\textit{Allowed standard prerequisites (editorial).}\begin{itemize}
\item The composite of two universal homeomorphisms is a universal homeomorphism, cited from the same chapter. Cited at proof lines 11441-11442. Disposes of the case in which the two factors are universal homeomorphisms.
\item A homeomorphism is an affine morphism, and an affine morphism is separated, both cited from the same chapter. Cited at proof lines 11479-11482. Shows that the auxiliary morphism in the last case is a closed immersion.
\item A section of a morphism of schemes is an immersion, cited from the schemes chapter. Cited at proof lines 11483-11485. Turns the identity relation between the auxiliary morphism and the homeomorphism into closedness of the auxiliary morphism.
\item The characterization of universal homeomorphisms, used here to conclude from bijectivity plus being a closed immersion. Cited at proof lines 11486-11488. Gives the final conclusion in the last case.
\end{itemize}\medskip
\ifdefined\hypertarget\hypertarget{sourcepointer-1}{}\fi
\noindent\textit{Source reference.} Original lemma 01S7: focused-seeds/AGP/sources/stacks-snapshot-20261011/tags.txt, source lines 2328--2328. Same-chapter reference used in the retained author text at source line(s) 1923, 11481. It is supplied only as a navigation pointer so that the retained text compiles; the counted item delivers the author's own statement and the author's complete proof, and no part of the referenced item is copied in or substituted for the counted argument.\par\medskip
\ifdefined\hypertarget\hypertarget{sourcepointer-2}{}\fi
\noindent\textit{Source reference.} Original lemma 0CEV: focused-seeds/AGP/sources/stacks-snapshot-20261011/tags.txt, source lines 15590--15590. Same-chapter reference used in the retained author text at source line(s) 11443, 12157. It is supplied only as a navigation pointer so that the retained text compiles; the counted item delivers the author's own statement and the author's complete proof, and no part of the referenced item is copied in or substituted for the counted argument.\par\medskip
\ifdefined\hypertarget\hypertarget{sourcepointer-3}{}\fi
\noindent\textit{Source reference.} Original lemma 04DE: focused-seeds/AGP/sources/stacks-snapshot-20261011/tags.txt, source lines 5724--5724. Same-chapter reference used in the retained author text at source line(s) 11372, 11416, 11480, 12247. It is supplied only as a navigation pointer so that the retained text compiles; the counted item delivers the author's own statement and the author's complete proof, and no part of the referenced item is copied in or substituted for the counted argument.\par\medskip
\ifdefined\hypertarget\hypertarget{sourcepointer-4}{}\fi
\noindent\textit{Source reference.} Original lemma 04DF: focused-seeds/AGP/sources/stacks-snapshot-20261011/tags.txt, source lines 5725--5725. Same-chapter reference used in the retained author text at source line(s) 11301, 11487, 11529, 11568, 11780, 11886, 14718, 14908. It is supplied only as a navigation pointer so that the retained text compiles; the counted item delivers the author's own statement and the author's complete proof, and no part of the referenced item is copied in or substituted for the counted argument.\par\medskip
\ifdefined\hypertarget\hypertarget{sourcepointer-5}{}\fi
\noindent\textit{Source reference.} Original lemma 01KT: focused-seeds/AGP/sources/stacks-snapshot-20261011/tags.txt, source lines 2062--2062. Same-chapter reference used in the retained author text at source line(s) 242, 1012, 2388, 6788, 11217, 11484, 12322, 12920. It is supplied only as a navigation pointer so that the retained text compiles; the counted item delivers the author's own statement and the author's complete proof, and no part of the referenced item is copied in or substituted for the counted argument.\par\medskip
\begin{lemma}
\label{lemma-permanence-universal-homeomorphism}
Let $f : X \to Z$ be the composition of two morphisms 
$g : X \to Y$ and $h : Y \to Z$.
If two of the morphisms $\{f, g, h\}$ are universal homeomorphisms,
so is the third morphism.
\end{lemma}

\begin{proof}
If both of $g$ and $h$ are universal homeomorphisms, so is $f$ by
Lemma \ref{lemma-composition-universal-homeomorphism}.

\medskip\noindent
Suppose both of $f$ and $g$ are universal homeomorphisms. We want to show
that $h$ is also. Now base change the diagram along an arbitrary morphism
$\alpha : Z' \to Z$ of schemes,  we get the following diagram with all
squares Cartesian:
$$
\xymatrix{
X' \ar[r]^{g'} \ar[d] & Y' \ar[r]^{h'} \ar[d] & Z' \ar[d] \\
X \ar[r]^{g} & Y \ar[r]^{h} & Z.
}
$$
Our assumption implies that the composition $f'= h' \circ g' : X' \to Z'$
and $g' : X' \to Y'$ are homeomorphisms, therefore so is $h'$.
This finishes the proof of $h$ being a universal homeomorphism.

\medskip\noindent
Finally, assume $f$ and $h$ are universal
homeomorphisms. We want to show that $g$ is a universal homeomorphism.
Let $\beta : Y' \to Y$ be an arbitrary morphism of schemes.
We get the following diagram with all squares Cartesian:
$$
\xymatrix{
X' \ar[r]^{g'} \ar[d] & Y' \ar[d]^{\gamma} & \\
X'' \ar[r]^{g''} \ar[d] & Y'' \ar[r]^{h''} \ar[d] & Y' \ar[d]^{h \circ \beta} \\
X \ar[r]^{g} & Y \ar[r]^{h} & Z.
}
$$
Here the morphism $\gamma : Y' \to Y''$ is defined by the universal property of
fiber products and the two morphisms
$id_{Y'} : Y' \to Y'$ and $\beta : Y' \to Y$.
We shall prove that $g'$ is a homeomorphism.
Since the property of being a homeomorphism has 2-out-of-3 property, we see
that $g''$ is a homeomorphism.
Staring at the top square, it suffices to prove that $\gamma$ is a universal
homeomorphism. Since $h''$ is a homeomorphism, we see that it is
an affine morphism by Lemma \ref{lemma-homeomorphism-affine}
and a fortiori separated (Lemma \ref{lemma-affine-separated}).
Since $h'' \circ \gamma$ is the identity, we see that $\gamma$
is a closed immersion by
Schemes, Lemma \ref{schemes-lemma-section-immersion}.
Since $h''$ is bijective, it follows that $\gamma$ is a bijective
closed immersion and hence a universal homeomorphism (for example by the
characterization in Lemma \ref{lemma-universal-homeomorphism}) as desired.
\end{proof}

\end{document}
